\section{Backbone and completion heads}
\label{sec:method}

\paragraph{Backbone.} Our backbone, GLWM-2, instantiates \cref{eq:local} with a pointwise embedding $E$,
the initial state $H=\mathrm{LN}$, a pointwise two-layer MLP decoder $D$, and the shared step
\begin{equation}
  h^{(t+1)} = \mathrm{LN}\Bigl(h^{(t)} + W_{\mathrm{out}}\,
  \mathrm{GELU}\bigl(\mathrm{Conv}^{(r)}_{3}\circ\cdots\circ\mathrm{GELU}\circ\mathrm{Conv}^{(1)}_{3}\,[h^{(t)}; e]\bigr)\Bigr).
  \label{eq:step}
\end{equation}
Each step applies $r$ kernel-3 convolutions to the current state concatenated with the input embedding
(input re-injection, or ``recall''), adds a residual connection, and normalises over channels at every
position. Each convolution widens the receptive field by one position on either side, so the growth rate
is exactly $r$, which we verify with the flip test. We use width $C=128$ and $r=1$. The backbone is
trained with the per-step cross-entropy averaged uniformly over $\Ttrain=20$ steps (the ``fixed prior''
objective) or with the progressive loss of \citet{bansal2022endtoend}; a pilot (\cref{sec:experiments})
selects the objective for each task.

\paragraph{Completion heads.} A completion head reads the latent state of each position and predicts
whether that position is \emph{done}, meaning that its prediction is correct now and will remain correct
at every later step. The target comes from the model's own trajectory up to a horizon $T$:
\begin{equation}
  z_i^{(t)} = \prod_{s=t}^{T}\ind\bigl[\yhat_i^{(s)} = y_i\bigr]
  = \min_{t\le s\le T}\ind\bigl[\yhat_i^{(s)} = y_i\bigr],
  \label{eq:done}
\end{equation}
that is, a reversed cumulative minimum of per-position correctness, and the head is trained with binary
cross-entropy against $z_i^{(t)}$. We study three variants. The \emph{post-hoc} head is a pointwise MLP
$p_i^{(t)} = \sigma\bigl(g(h_i^{(t)})\bigr)$ fitted after the backbone has been frozen, on rollouts of
length $2\Ttrain$. The \emph{joint} head has the same form but is trained alongside the backbone with an
auxiliary binary cross-entropy of weight 1. Unlike ACT and PonderNet, it never reweights the per-step
task loss: the backbone receives the same task gradient at every step, so predicting completion is a
separate problem rather than a halting distribution laid over the task loss. Finally, the
\emph{propagating} head also reads its neighbours' previous done probabilities,
\begin{equation}
  p_i^{(t)} = \sigma\Bigl(g\bigl(h_i^{(t)}\bigr) + \mathrm{Conv}_3\bigl[p^{(t-1)}, \mathbf 1\bigr]_i\Bigr),
  \label{eq:prop}
\end{equation}
which mirrors the pending-bit automaton of \cref{prop:automaton}: a position can learn to wait until
its left neighbour reports that it is done.

\paragraph{Halting rule.} An instance halts at the first step at which every valid position is done with
probability above a threshold $\theta$, or at the budget $\Tmax$:
\begin{equation}
  \tau_\theta(x) = \min\Bigl(\bigl\{\,t\ge 1 : \min_{0\le i<P} p_i^{(t)}(x) > \theta\,\bigr\}\cup\{\Tmax\}\Bigr),
  \label{eq:halt}
\end{equation}
and returns $\yhat^{(\tau_\theta(x))}(x)$ (\cref{alg:halt} in \cref{app:details}). We select $\theta$ on
ID-val as the cheapest threshold whose accuracy is within 0.5 points of the best. Because the rule takes a
minimum over positions, its cost is set by the slowest position rather than by the input length.

\paragraph{Baselines.} Every baseline uses the same backbone and the same selection protocol. As
references we report a fixed step count $T$ selected on ID-val, the oracle lower bound $\LB(x)$, the
any-step oracle (the first step with the right answer) and the stable oracle (the first step from which
the answer stays right). Writing $\pi_i^{(t)} = \operatorname{softmax}D(h_i^{(t)})$ for the predictive
distribution at position $i$, four post-hoc rules halt once a scalar statistic crosses a threshold:
\begin{equation}
  \begin{aligned}
    S_{\mathrm{msp}}^{(t)} &= \min_i \max_\gamma \pi_{i\gamma}^{(t)}, &\quad
    S_{\mathrm{ent}}^{(t)} &= \max_i \mathrm{H}\bigl(\pi_i^{(t)}\bigr),\\
    S_{\mathrm{KL}}^{(t)} &= \max_i \mathrm{KL}\bigl(\pi_i^{(t)}\,\big\|\,\pi_i^{(t-1)}\bigr), &\quad
    S_{\Delta h}^{(t)} &= \frac{\lVert h^{(t)}-h^{(t-1)}\rVert_F}{\lVert h^{(t)}\rVert_F},
  \end{aligned}
  \label{eq:posthoc}
\end{equation}
where the first halts when it rises above its threshold and the other three when they fall below it.
Note that max-softmax and entropy are per-position statistics combined by a worst case, whereas the
hidden-state change pools over the whole input. We further include $k$-step prediction stability
\citep{zhou2020pabee}, maximum confidence over steps \citep{fan2025looped}, and a pooled correctness head
in the style of TRM \citep{jolicoeur2025trm} fitted post hoc. The jointly
trained rules are ACT \citep{graves2016act}, PonderNet with its KL term towards a geometric prior
\citep{banino2021pondernet}, a joint pooled correctness head, and the pooled gate of our earlier GLWM
model (v1) with its initialisation fixed.
